% Timeless material: lines picked by what happened, reactions to events, and the cards.

\section{Blue reacts: the first coin moves}\label{react:first-move}
\on{first-move}

Blue (flow): Hey!
\todo{draft}

% after the first coin moves.

\section{Blue reacts: down to eleven}\label{react:blue-reaches-eleven}
\on{blue-reaches-eleven}

Blue (flow): I liked that one.
\todo{draft}

% when Blue is down to 11.

\section{They react: Red passes eight}\label{react:red-above-eight}
\on{red-above-eight}

Blue (flow): Whoa. Equal, not upside down.
\todo{draft}

Red (flow): Tempting. But give him one back.
\todo{draft}

% when Red goes above 8.

\section{The bet: a coffee}\label{bet:coffee}
\when{bet=coffee}

Blue (aside, flow): A coffee. Careful with your money. I respect that.
\todo{draft}

\section{The bet: a lunch}\label{bet:lunch}
\when{bet=lunch}

Red (aside, flow): A lunch. Fair. The winner picks the place.
\todo{draft}

\section{The bet: a vacation}\label{bet:vacation}
\when{bet=vacation}

Blue (aside, flow): A vacation! Now it's serious.
\todo{draft}

\section{The bet: one percent}\label{bet:percent}
\when{bet=percent}

Red (aside, flow): One percent. Careful. That's how it starts.
\todo{draft}

\section{After the run: someone else wins}\label{run:other}
\when{winner=other}

Blue (flow): Red? Where are you?

Red (flow): Here. What do you have?

Blue (flow): \val{blue}. You?
\todo{draft}

Red (flow): \val{red}.
\todo{draft}

Red (flow): Told you.

% when someone else finishes richest (the most likely case); {blue} and {red}
% are what they actually hold, in dollars (everyone started with $100).

\section{After the run: Blue wins}\label{run:blue}
\when{winner=blue}

Blue (flow): See? Talent rises.
\todo{draft}

Red (flow): Run it again and see if your talent comes along.
\todo{draft}

% when Blue finishes richest.

\section{After the run: Red wins}\label{run:red}
\when{winner=red}

Red (flow): Look at that. Am I a genius now?
\todo{draft}

Blue (flow): …Run it again.
\todo{draft}

% when Red finishes richest.

\section{Why they won: after one run}\label{why:once}
\when{runs=one}

Red: The paper found the reason after the result.
The coin never saw a thing.
\todo{draft}

\section{Why they won: after several runs}\label{why:again}
\when{runs=several}

Red: New winner. New reason.
The paper never prints "chance."
\todo{draft}

% instead of [why.after] once the reader has run the room more than once.

% Red's one line about the reader's last room, whichever fits.

\section{The dial: zero}\label{dial:zero}
\when{stake=zero}

Red (flow): Zero: nothing moves.
That's no game at all.
\todo{draft}

\section{The dial: slower}\label{dial:slow}
\when{stake=slow}

Red (flow): \val{stake}: slower.
The richest holds \val{share}.
\todo{draft}

\section{The dial: the stake you watched}\label{dial:same}
\when{stake=same}

Red (flow): \val{stake}: the stake you watched.
The richest holds \val{share}.
\todo{draft}

\section{The dial: faster}\label{dial:fast}
\when{stake=fast}

Red (flow): \val{stake}: faster.
The richest holds \val{share}.
\todo{draft}

\section{The dial: everything}\label{dial:all}
\when{stake=all}

Red (flow): Everything on the table:
one loss and you're out. The richest holds \val{share}.
\todo{draft}

\section{The tax game: won}\label{game:won}
\when{game=won}

Red: Thirty seconds, and still \val{count} players.
Your hand did it.
\todo{draft}

\section{The tax game: lost}\label{game:lost}
\when{game=lost}

Red: It closed after \val{seconds} seconds.
\val{taps} taps weren't enough.
\todo{draft}

\section{The cards}\label{cards}

\begin{description}
  \item[card.stack] Cards
  \item[card.play] Play with it
  \item[card.close] Close
\end{description}

\section{Card: The rule}\label{card:rule}

\begin{description}
  \item[card.rule.title] The rule
  \item[card.rule.1] Pick two people at random.
  \item[card.rule.2] Each puts in \val{stake} of what the poorer one has.
  \item[card.rule.3] Flip a coin. The winner takes both.
  \item[card.rule.4] Do it again.
\end{description}

% the run plays at exactly a quarter of the poorer fortune — the same bet Red
% names (owner 2026-09-26: half was too jittery; it used to be 35%, then
% half).

% the first concept card (brief 1.6): like the rule card in a board game, it
% shows the rule that is running right now, read from the room's own settings
% — {stake} is the live stake. It drops into the stack in the corner once Red
% has said the whole rule (Scene 10), and opens at "Remember the rule."

\section{Card: Histogram}\label{card:histogram}

\begin{description}
  \item[card.histogram.title] Histogram
  \item[card.histogram.1] Sort everyone into piles by how much they have.
  \item[card.histogram.2] A pile's height is how many people are in it.
  \item[card.histogram.3] A multiplying ruler shows what an ordinary one squeezes into a corner.
\end{description}

\section{Card: Gini}\label{card:gini}

\begin{description}
  \item[card.gini.title] Gini
  \item[card.gini.1] Line everyone up, poorest first, and add up their money as you go.
  \item[card.gini.2] 0: everyone has the same. 1: one owner, nothing for the rest.
\end{description}

\section{Card: Effective participants}\label{card:participants}

\begin{description}
  \item[card.participants.title] Effective participants
  \item[card.participants.1] How many equal fortunes would be just as concentrated as this room.
  \item[card.participants.2] Everyone equal: all of them. One owner: 1. No money: you don't count.
  \item[participation.formula] Open the black box: write each person's share of the room as sᵢ. Square the shares, add them, then take the reciprocal: 1 / Σsᵢ². This is the inverse Herfindahl concentration index.
\end{description}

% a claim note (the fold-out "Open the black box"), not dialogue — on the
% card.

\section{Card: Turnover}\label{card:turnover}

\begin{description}
  \item[card.turnover.title] Turnover
  \item[card.turnover.1] How much of all the money changes hands in one round: one trade per person.
  \item[card.turnover.2] A busy room moves a lot of its money. A concentrated one barely moves any, however many trades it makes.
\end{description}

\section{Card: Where it ends}\label{card:limit}

\begin{description}
  \item[card.limit.title] Where it ends
  \item[card.limit.1] Keep playing forever: with probability one, one person ends up with everything.
  \item[card.limit.2] A limit, not a date. At any moment the room can still wobble.
  \item[card.limit.3] Proof: Börgers \& Greengard, 2023.
\end{description}

\section{Card: The stake}\label{card:stake}

\begin{description}
  \item[card.stake.title] The stake
  \item[card.stake.1] How much of the poorer one's money goes on each toss.
  \item[card.stake.2] Any stake above zero ends the same way. Smaller only takes longer. Zero is no game.
\end{description}

\section{Card: The shared levy}\label{card:levy}

\begin{description}
  \item[card.levy.title] The shared levy
  \item[card.levy.1] Once a round: the same share of every fortune into one pool.
  \item[card.levy.2] The pool comes back in equal parts.
  \item[card.levy.3] Below the average gains, the average breaks even, above it pays in. Nothing is lost.
\end{description}
